% TOP_PROBLEM: {"id": 20000924, "title": "The minimal empty-core tight path and the 2026 uniformity-four resolution", "category_id": 2, "category": {"id": 2, "name": "combinatorics", "display_name": "Combinatorics", "description": "Counting problems, graph theory, discrete structures.", "slug": "combinatorics", "order_index": 2, "created_at": "2026-07-31T15:26:25.671Z"}, "status": "partially_solved", "problem_number": "AIM-COMBINATORICS-0049", "set_id": 14, "statusQualification": "The all-uniformities workshop assertion is retained; Nie proves the relevant long-path cases in uniformities3 and4, so the remaining general range is k>=5."}
% FIRST_POSTED: 2026-10-01
% ENTRY_KIND: statement_only
% UnsolvedMath ID 20000924; source catalogue code AIM-COMBINATORICS-0049
% !TeX program = lualatex
% Definition and sources only; this file is not an LLM solution attempt.
\documentclass[11pt,a4paper]{article}
\usepackage[margin=25mm]{geometry}
\usepackage{fontspec}
\setmainfont{lmroman10-regular.otf}[BoldFont=lmroman10-bold.otf,
 ItalicFont=lmroman10-italic.otf,BoldItalicFont=lmroman10-bolditalic.otf]
\usepackage{amsmath,amssymb,mathtools}
\usepackage{unicode-math}
\setmathfont{latinmodern-math.otf}
\AtBeginDocument{\let\setminus\smallsetminus}
\usepackage{xurl,hyperref}
\hypersetup{colorlinks=true,urlcolor=blue}
\setcounter{secnumdepth}{0}
\setlength{\parindent}{0pt}
\setlength{\parskip}{5pt}
\setlength{\emergencystretch}{3em}
\begin{document}
\section*{The minimal empty-core tight path and the 2026 uniformity-four resolution}\label{20000924}
{\small UnsolvedMath ID 20000924 · AIM-COMBINATORICS-0049 · Combinatorics · AIM Workshop Problem Lists}
\subsection{Definitions and mathematical statement}
For integers $k\ge2$ and $s\ge1$, the tight $k$-uniform path with
$s$ edges, denoted $P_s^{(k)}$, has vertex set $[s+k-1]$ and edges
\[
             \{i,i+1,\ldots,i+k-1\},\qquad 1\le i\le s.
\]
The number $s$ counts edges, not vertices. Write $K_t^{(k)}$ for the
complete $k$-uniform hypergraph on $t$ vertices. For $t\ge k$, let
$r_k(P_s^{(k)},K_t^{(k)})$ be the least $N$ such that every coloring
of the $k$-element subsets of $[N]$ red or blue contains a red copy of
$P_s^{(k)}$ or a blue copy of $K_t^{(k)}$. All copies are injective
and ordinary, rather than induced, hypergraph copies.

The question is whether for every fixed $k\ge2$ there exists an integer
$s_0=s_0(k)$ such that for every fixed $s\ge s_0$,
\[
               r_k(P_s^{(k)},K_t^{(k)})=\Theta_{k,s}(t^{k-1})
                       \quad(t\to\infty).
\]
Explicitly, positive constants $c_{k,s},C_{k,s}$ and a threshold
$t_0(k,s)$ must give both corresponding inequalities for all
$t\ge t_0(k,s)$. The exponent is prescribed, while the multiplicative
constants may depend on the path. Neither $k$ nor $s$ grows with $t$
in this assertion. The central issue is whether making the fixed path
sufficiently long forces the full power $t^{k-1}$ in every uniformity.
Nie's Theorem 1.1 establishes this order for nontrivial tight trees in
uniformities three and four, hence for the corresponding sufficiently
long paths. His Conjecture 1 explicitly retains uniformities at least
five; these are the remaining range of the general question here.
\subsection{Short English statement}
For every uniformity at least five, do sufficiently long fixed tight paths have clique Ramsey numbers of order the expected full power of the clique size?
\subsection{Sources}
\begin{itemize}
\item[S1] ULAM.AI and the UnsolvedMath Contributors, supplied problems.json, record 20000924 (AIM-COMBINATORICS-0049), AIM Workshop Problem Lists. Catalogue identity and source transcription; CC BY 4.0.
\url{https://huggingface.co/datasets/ulamai/UnsolvedMath}.
\item[S2] American Institute of Mathematics, Graph Ramsey theory, item 2. Original workshop question underlying AIM-COMBINATORICS-0049.
\url{https://www.overleaf.com/read/mnvcscjjysvg}.
\item[S3] Jiaxi Nie, On tight tree-complete hypergraph Ramsey numbers, Theorem 1.1 and Conjecture 1. Distinguishes proved uniformities three and four from the higher-uniformity target.
\url{https://arxiv.org/pdf/2412.19461}.
\end{itemize}
\end{document}
